\documentclass[border=12pt]{standalone}
\usepackage[utf8]{inputenc}
\usepackage[T1]{fontenc}
\usepackage[spanish,es-noshorthands]{babel}
\usepackage{tikz}
\usetikzlibrary{positioning, shapes.geometric, shapes.symbols, arrows.meta, calc, fit, backgrounds}

% ---- Paleta de colores ----
\definecolor{azulPC}{RGB}{33,102,172}      % PC / App Qt
\definecolor{azulPCbg}{RGB}{222,235,247}
\definecolor{verdeESP}{RGB}{35,132,67}     % ESP32
\definecolor{verdeESPbg}{RGB}{225,245,229}
\definecolor{grisRed}{RGB}{110,110,110}    % Red WiFi
\definecolor{grisRedbg}{RGB}{238,238,238}
\definecolor{naranjaHW}{RGB}{217,120,20}   % Hardware
\definecolor{naranjaHWbg}{RGB}{253,237,214}
\definecolor{textoOsc}{RGB}{40,40,40}

\begin{document}
\begin{tikzpicture}[
  font=\sffamily,
  >={Latex[length=2.6mm]},
  % --- estilos de bloque ---
  bloque/.style={
    rectangle, rounded corners=6pt, draw, line width=1pt,
    align=center, inner sep=8pt, minimum height=22mm, text=textoOsc
  },
  pc/.style={bloque, draw=azulPC, fill=azulPCbg, minimum width=46mm},
  esp/.style={bloque, draw=verdeESP, fill=verdeESPbg, minimum width=46mm},
  hw/.style={
    rectangle, rounded corners=4pt, draw=naranjaHW, fill=naranjaHWbg,
    line width=0.9pt, align=left, inner sep=6pt, text=textoOsc,
    minimum width=52mm
  },
  nube/.style={
    cloud, cloud puffs=13, cloud puff arc=110, aspect=2.3,
    draw=grisRed, fill=grisRedbg, line width=1pt, align=center,
    inner sep=2pt, minimum width=52mm, minimum height=30mm, text=textoOsc
  },
  etiq/.style={font=\sffamily\footnotesize, align=center, text=textoOsc},
  etiqRet/.style={font=\sffamily\footnotesize, align=center, text=grisRed}
]

% ================= BLOQUES PRINCIPALES (fila horizontal) =================
\node[pc]  (pc)  {\textbf{PC — App Qt}\\\textbf{(Control del LED)}\\[2pt]
  \footnotesize Botones \textbf{ENCENDER} / \textbf{APAGAR}\\
  \footnotesize$\rightarrow$ envían mensaje\\
  \footnotesize WebSocket (JSON)\\[2pt]
  \scriptsize(también sirve un navegador)};

\node[nube] (red) [right=34mm of pc] {\textbf{Red WiFi}\\\textbf{SSID GWN571D04}};

\node[esp]  (esp) [right=34mm of red] {\textbf{ESP32-S3 (Arduino)}\\\textbf{Servidor WebSocket}\\[2pt]
  \footnotesize Escucha en\\
  \footnotesize\texttt{ws://<IP>:81}};

% ================= FLECHAS PC <-> RED <-> ESP32 =================
% Ida: PC -> Red
\draw[->, line width=1.1pt, azulPC] (pc.east) -- node[etiq, above, yshift=1pt]{WebSocket \texttt{ws://<IP>:81}}
      node[etiq, below, yshift=-1pt]{\scriptsize\texttt{\{"tipo":"led","encendido":true|false\}}} (red.west);
% Ida: Red -> ESP32
\draw[->, line width=1.1pt, verdeESP] (red.east) -- node[etiq, above, yshift=1pt]{mensajes JSON}
      node[etiq, below, yshift=-1pt]{\scriptsize\texttt{\{"tipo":"led", ...\}}} (esp.west);

% Retorno: ESP32 -> Red -> PC (arco inferior)
\draw[->, line width=1.0pt, grisRed]
   (esp.south west) ++(0,-2mm)
   to[out=210, in=330] node[etiqRet, below, yshift=-1pt]{\scriptsize\texttt{\{"tipo":"estado\_led","encendido":..\}}}
   ($(red.south east)+(0,-2mm)$);
\draw[->, line width=1.0pt, grisRed]
   ($(red.south west)+(0,-2mm)$)
   to[out=210, in=330] node[etiqRet, below, yshift=-1pt]{\scriptsize\texttt{\{"tipo":"saludo"\}} al conectar}
   ($(pc.south east)+(0,-2mm)$);

% ================= BLOQUE HARDWARE (colgando del ESP32) =================
\node[hw] (hw) [below=22mm of esp] {%
  \textbf{Hardware}\\[3pt]
  \textbullet~LED \texttt{GPIO4} + R 330\,$\Omega$\\
  \textbullet~LED RGB \texttt{GPIO48}\\
  \textbullet~Pulsador \texttt{GPIO15} (se muestra su estado)};

\draw[->, line width=1.1pt, naranjaHW] (esp.south) -- node[etiq, right, xshift=1pt, text=naranjaHW]{control / lectura} (hw.north);

% ================= NOTA AL PIE =================
\node[etiq, below=8mm of hw, text width=175mm] (nota) {%
  \itshape Mismo ida y vuelta que el TC4 (WebSocket + JSON con campo \texttt{tipo}); en el TC4 se agregan más comandos.};

\end{tikzpicture}
\end{document}
